\section{Q1: Surrogate Modeling and Low-Tail Screening}
\label{sec:q1-surrogate}

\subsection{Objective and violation score}

Question~1 asks for a composition-based surrogate that ranks candidate
perovskites by their violation of five material requirements.  For a
composition $z=(A,B,X)$, let
\begin{equation}
  p(d)=\left(\frac{\max(0,d)}{0.5}\right)^2 .
  \label{eq:q1-hinge}
\end{equation}
Four predicted properties enter the score.  We denote formation energy and
band gap by $e_{\mathrm{form}}$ and $E_g$, respectively.  The valence- and
conduction-band edges are $E_{\mathrm{VBM}}$ and $E_{\mathrm{CBM}}$.  The
official violation score is
\begin{align}
  V={}&p(e_{\mathrm{form}}-0.21)
      +p(1.4-E_g)+p(E_g-3.1) \notag\\
     &+p(5.73-E_{\mathrm{VBM}})
      +p(E_{\mathrm{CBM}}-4.50).
  \label{eq:q1-violation}
\end{align}
All energies are measured in eV \cite{apmcm2026problemD}.  A zero score
allows equality at a threshold, whereas strict feasibility uses open
inequalities.  We therefore model the continuous score in
Equation~\eqref{eq:q1-violation} and keep strict feasibility as a separate
concept.

The labeled training and validation files contain 11,256 and 3,969 rows,
respectively.  Neither split contains a zero-violation or strictly feasible
record.  Moreover, $96.16\%$ of the training rows and $96.45\%$ of the
validation rows have zero band gap.  Since $p(1.4)=7.84$, a metallic
prediction alone creates a substantial penalty.  These facts make global
regression error insufficient for model selection: the surrogate must also
preserve the ordering of the small low-$V$ subset.  File~04 contains 3,703
unlabeled compositions and is used only for blind inference; its true error
and target recall are not observable.

\subsection{Composition features and candidate surrogates}

The permitted input is the ordered categorical triple $(A,B,X)$.  The two
cation sites remain distinct, candidates with $A=B$ are retained, and the
identifier, row order, pool membership, and physical labels are excluded
from the feature vector.  We compare three levels of modeling complexity.

First, direct controls predict $V$ itself: a training-set median, one-hot
Ridge regression, a random forest, and histogram gradient boosting (HGB).
An additional Ridge ablation includes deterministic $A{:}X$ and $B{:}X$
interactions and an $A=B$ indicator.  Second, property-based models predict
the physical quantities and insert them into
Equation~\eqref{eq:q1-violation}.  Third, two-stage models treat the highly
imbalanced band-gap term separately.  A tail-weighted direct control uses
the training weight
\begin{equation}
  w(V)=1+4\exp(-V/5),
  \label{eq:q1-tail-weight}
\end{equation}
so that low-score records receive more influence without changing the
target definition.

All HGB comparisons use the same configuration: learning rate $0.075$,
180 boosting iterations, at most 31 leaves, a minimum of 20 samples per
leaf, and $\ell_2$ regularization of $1.0$.  Early stopping is disabled and
the categorical encoders are fitted separately inside every training
partition.  This matched design isolates differences in the target
construction rather than changes in estimator capacity.

\subsection{Two-stage band-gap penalty mixture}

The selected model uses three regression heads for
$e_{\mathrm{form}}$, $E_{\mathrm{VBM}}$, and $E_{\mathrm{CBM}}$.  Define the
combined gap penalty
\begin{equation}
  G(E_g)=p(1.4-E_g)+p(E_g-3.1).
  \label{eq:q1-gap-penalty}
\end{equation}
A classifier estimates
$q(z)=\Pr(E_g>0\mid z)$, while a regression head fitted only to positive-gap
records estimates $\widehat E_g^{+}(z)$.  Negative outputs of this head are
floored at zero.  Instead of applying a hard class decision, the point
penalty mixture is
\begin{equation}
  \widehat G_{\mathrm{mix}}(z)
   =[1-q(z)]G(0)+q(z)G\!\left(\max(0,\widehat E_g^{+}(z))\right).
  \label{eq:q1-gap-mixture}
\end{equation}
The final surrogate is therefore
\begin{align}
  \widehat V(z)={}&p(\widehat e_{\mathrm{form}}-0.21)
  +p(5.73-\widehat E_{\mathrm{VBM}})
  +p(\widehat E_{\mathrm{CBM}}-4.50) \notag\\
  &+\widehat G_{\mathrm{mix}}(z).
  \label{eq:q1-surrogate}
\end{align}

The hard-gate control assigns a positive gap only when $q\geq0.5$, whereas
the mean-gap control evaluates $G(q\widehat E_g^{+})$.  On file~03 the
positive-gap classifier had average precision $0.5774$, precision $0.8500$,
and recall $0.2411$ at the $0.5$ threshold.  It would classify 49 of the 70
observed $V<5$ rows as zero-gap cases.  The conditional positive-gap
regressor had an MAE of $0.8290$ eV over 141 positive-gap rows.  These
diagnostics motivated retaining soft alternatives to the hard gate; they do
not establish a causal explanation for their ranking performance.

\subsection{Grouped training, validation, and selection}

Development uses five folds on file~02.  Its 1,608 ordered $(A,B)$ pairs are
kept intact, so all seven $X$ variants of a pair enter the same fold.  Every
candidate model is evaluated on the same partitions.  Besides MAE, RMSE,
and Spearman correlation, we report MAE on the observed $V<5$ subset and
ranking-based screening metrics.  Here $V<5$ is a project-defined operational
low-tail threshold chosen to yield a rare but nonempty subset in both labeled
files; it is not an official feasibility condition.  If $T_K$ and
$\widehat T_K$ denote the
true and predicted lowest-$K$ sets, respectively, the overlap is
\begin{equation}
  R_K=\frac{|T_K\cap\widehat T_K|}{K},\qquad
  S=\frac{R_{25}+R_{50}}{2}.
  \label{eq:q1-screening-score}
\end{equation}
Ties are resolved by numerical material identifier.  The screening score
$S$ is used jointly with global and low-tail errors, rather than as a claim
that any selected material is physically feasible.

Four configurations were frozen before the formal file~03 comparison:
direct HGB, tail-weighted direct HGB, the mean-gap property model, and the
penalty mixture.  The global-MAE acceptance limit was fixed at $1.05$ times
the best shortlisted direct-model MAE, giving $2.8432$.  Within a
two-percentage-point screening-score tie band, low-$V$ MAE was the next
criterion.  Once the specification was selected, the Q1 release model
\texttt{q1\_final} was fitted only on file~02 and applied to file~04.  A
separate \texttt{q23\_final} model uses the same frozen specification but is
refitted on files~02 and~03 for the later candidate pools; its predictions
are not used as Q1 validation evidence.

\subsection{Grouped ablation results}

Across the principal five-fold comparisons, the mixture obtained the largest mean screening score,
$39.8\%$, while maintaining an MAE of $2.649$.  Direct HGB reached
$27.0\%$ screening with MAE $2.877$.  On matched folds, the mixture's
screening score exceeded direct HGB by an average of 12.8 percentage points
(paired standard deviation 9.1 points) and the ordinary-gap property model
by 6.4 points (standard deviation 3.5 points).  Its 0.8-point mean advantage
over the mean-gap model changed sign across folds, so this small difference
is not evidence of uniform superiority.  The ordinary-gap and mean-gap
models also achieved smaller $V<5$ MAE than the selected mixture, exposing
the trade-off between pointwise low-tail error and Top-$K$ recovery.

\subsection{File~03 validation and low-tail screening}

The frozen comparison is documented in Appendix
Table~\ref{tab:q1-validation}.  The
tail-weighted direct model had the largest screening score, $38\%$, but its
MAE of $2.8655$ exceeded the predeclared $2.8432$ limit.  The penalty
mixture passed the limit and achieved the lowest MAE and RMSE among the four
frozen candidates, together with a screening score of $36\%$.  It was
therefore selected without revising the rule after observing the results.

At budgets $K=10,25,50,100$, the selected model recovered respectively
1, 8, 20, and 41 members of the true lowest-$K$ sets.  The minimum observed
score among its selected candidates was $0.0196$ at every budget.  Relative
to direct HGB, the selected model was associated with an $8.7\%$ lower MAE,
twice as many Top-25 hits, and an increase from 13 to 20 Top-50 hits.  The
number of true zero-gap records in the predicted Top-50 also fell from 29
to 13.  These are matched empirical comparisons, not causal effects of the
mixture construction.

The archived five-component diagnostic shows that, across all validation
rows, the observed and predicted
mean components are broadly aligned.  In the 70-row $V<5$ subset, however,
the predicted lower-gap and CBM penalties are much larger than their
observed means.  This pattern is consistent with the remaining low-tail
MAE and cautions against interpreting fine differences within the predicted
minimum as physical certainty.

The expanded post-freeze panel also evaluated a model that directly fits
the conditional expected positive-gap penalty.  On file~03 it obtained MAE
$2.4624$ and screening score $41\%$, compared with $2.4715$ and $36\%$ for
the frozen mixture.  In the original grouped development folds, however,
its screening score was $35.2\%$, below the mixture's $39.8\%$.  Because
these additional file~03 comparisons were performed after selection, they
serve as sensitivity evidence and do not justify replacing the frozen
release.

\subsection{Stability and finite-budget screening value}

With the original ordered-pair folds fixed, changing the HGB estimator seed
from 42 to 1729 or 31415 left the compared metrics unchanged.  With the
estimator seed fixed at 42 and two alternative group allocations, the
mixture achieved screening scores of $40.2\%$ and $42.6\%$; direct HGB
obtained $28.0\%$ and $26.0\%$.  The unchanged fixed-fold seed results provide
limited information because this HGB configuration is effectively deterministic
under those inputs.  The two alternative group allocations show partition
sensitivity over the tested splits, not a distribution-free stability guarantee.

To measure value under a finite verification budget, we replayed the frozen
file~03 ranking against uniform sampling without replacement.  For a pool
of size $N$ containing $H$ records with $V<5$, a random budget of $K$ has
expected low-tail hits
\begin{equation}
  \mathrm{E}[Y]=\frac{KH}{N}.
  \label{eq:q1-random-hits}
\end{equation}
sampling.  Appendix Table~\ref{tab:q1-budget} shows that the surrogate found the
validation-set minimum at every tested budget and concentrated many more
$V<5$ records than the random expectation.

For $K=50$, the exact probability that uniform sampling attains a minimum no
larger than $0.0196$ is $1.26\%$, while the exact probability of at least 23
low-tail hits is $6.64\times10^{-29}$.  These within-pool probabilities do not
correct for file~03 having participated in model selection.

The same retrospective calculation was repeated under the Q2-like
$X=\mathrm{NO}_2$ subset and the Q3-like
$X\in\{\mathrm{O}_3,\mathrm{NO}_2,\mathrm{FO}_2\}$ subset of file~03.  At
$K=50$, the proxy found 9 and 14 low-$V$ records, respectively, compared
with random expectations of $0.9700$ and $1.2346$; the corresponding
enrichment factors were 9.28 and 11.34.  These subpools overlap and file~03
was used in model selection, so the replay quantifies screening value on
this observed validation set only.  It does not estimate success
probabilities for file~04 or for unseen chemical domains.

\subsection{Strengths, limitations, and Q1 conclusion}

The principal strength of the approach is its auditable separation of data
roles.  Ordered-pair grouping prevents the seven $X$ variants of a pair from
crossing development folds; the frozen selection rule combines global
accuracy with low-tail retrieval; and the probabilistic gap mixture avoids
making the low-recall hard classifier the sole determinant of the gap
penalty.  Versioned model files and prediction tables also separate the
Q1-only fit from the later Q2/Q3 scoring fit.  Within the tested data, the
results support using the surrogate as a finite-budget ranking device.

Several limitations remain.  First, neither labeled split contains a
zero-violation or strictly feasible record, so recall of the rare official
targets cannot be measured.  Second, file~03 served both model selection
and the reported formal comparison; the later expanded panel and random
budget replay are therefore descriptive post-freeze evidence rather than a
new independent validation set.  Third, the positive-gap probabilities
were not separately calibrated.  Fourth, the VBM and CBM MAEs rise from
$0.061$ and $0.061$ eV globally to $0.667$ and $0.664$ eV within the
observed low-$V$ subset, limiting confidence in fine ordering among the
best-looking candidates.  Finally, the model uses composition categories
only and provides no guarantee for unfamiliar elements or for physical
feasibility beyond the supplied labels.

Accordingly, the 3,703 file~04 predictions are blind surrogate scores, not
verified discoveries.  Their holdout MAE, target recall, and strict
feasibility rate are unknown.  Machine-level reproduction and artifact
checks have passed, but independent team sign-off and the official
platform-field check remain pending.  The Q1 evidence therefore supports
low-tail prioritization and budgeted follow-up, while any selected material
still requires independent physical or computational verification.
